%%%%%%%%%%%%%%%%%%%%%%%%%%%%%%%%%%%%%%%%%%%%%%%%%%%%%%%%%%%%%%%%%
% PRESENTACIÓN: INTELIGENCIA ARTIFICIAL APLICADA PARA UNIVERSITARIOS
% Autor: Edison Achalma
% Institución: Universidad Nacional de San Cristóbal de Huamanga
% Duración: 2 horas (120 minutos)
% Fecha de generación: 2025
%%%%%%%%%%%%%%%%%%%%%%%%%%%%%%%%%%%%%%%%%%%%%%%%%%%%%%%%%%%%%%%%%

\documentclass[aspectratio=169,xcolor={dvipsnames}]{beamer}
% aspectratio=169 para formato 16:9 (moderno, ideal para proyectores HD)

%%%%%%%%%%%%%%%%%%%%%%%%%%%%%%%%%%%%%%%%%%%%%%%%%%%%%%%%%%%%%%%%%
% PAQUETES ESENCIALES
%%%%%%%%%%%%%%%%%%%%%%%%%%%%%%%%%%%%%%%%%%%%%%%%%%%%%%%%%%%%%%%%%

% Codificación y idioma
\usepackage[utf8]{inputenc}
\usepackage[spanish,es-tabla]{babel} % es-tabla para "Tabla" en lugar de "Cuadro"
\usepackage[T1]{fontenc}

% Tema y apariencia
\usetheme{Madrid}
\usecolortheme{default}

% Gráficos y colores
\usepackage{graphicx}
\usepackage{xcolor}
\usepackage{tikz}
\usetikzlibrary{shapes,arrows,positioning,shadows,fit}

% Tablas profesionales
\usepackage{booktabs}
\usepackage{array}
\usepackage{multirow}

% Hipervínculos
\usepackage{hyperref}
\hypersetup{
	colorlinks=true,
	linkcolor=white,      % Links en TOC blancos
	urlcolor=cyan,        % URLs en celeste
	citecolor=blue
}

% Íconos y símbolos
\usepackage{fontawesome5} % Requiere instalación: tlmgr install fontawesome5

% Listas mejoradas
\usepackage{enumitem}

% Código (para slides de herramientas de programación)
\usepackage{listings}
\usepackage{xcolor}

%%%%%%%%%%%%%%%%%%%%%%%%%%%%%%%%%%%%%%%%%%%%%%%%%%%%%%%%%%%%%%%%%
% CONFIGURACIONES DE DISEÑO PERSONALIZADAS
%%%%%%%%%%%%%%%%%%%%%%%%%%%%%%%%%%%%%%%%%%%%%%%%%%%%%%%%%%%%%%%%%

% Colores institucionales UNSCH (personalizar según identidad visual)
\definecolor{azulUNSCH}{RGB}{0,51,102}      % Azul institucional
\definecolor{naranjaUNSCH}{RGB}{255,102,0}  % Naranja acento
\definecolor{grisClaro}{RGB}{240,240,240}   % Gris fondos

% Aplicar colores al tema
\setbeamercolor{structure}{fg=azulUNSCH}
\setbeamercolor{title}{fg=white,bg=azulUNSCH}
\setbeamercolor{frametitle}{fg=white,bg=azulUNSCH}
\setbeamercolor{block title}{fg=white,bg=azulUNSCH}
\setbeamercolor{block body}{fg=black,bg=grisClaro}
\setbeamercolor{alerted text}{fg=naranjaUNSCH}

% Eliminar símbolos de navegación (limpia la presentación)
\setbeamertemplate{navigation symbols}{}

% Configurar footer con número de slide
\setbeamertemplate{footline}{
	\leavevmode%
	\hbox{%
		\begin{beamercolorbox}[wd=.333333\paperwidth,ht=2.25ex,dp=1ex,center]{author in head/foot}%
			\usebeamerfont{author in head/foot}\insertshortauthor
		\end{beamercolorbox}%
		\begin{beamercolorbox}[wd=.333333\paperwidth,ht=2.25ex,dp=1ex,center]{title in head/foot}%
			\usebeamerfont{title in head/foot}\insertshorttitle
		\end{beamercolorbox}%
		\begin{beamercolorbox}[wd=.333333\paperwidth,ht=2.25ex,dp=1ex,right]{date in head/foot}%
			\usebeamerfont{date in head/foot}\insertshortdate{}\hspace*{2em}
			\insertframenumber{} / \inserttotalframenumber\hspace*{2ex}
		\end{beamercolorbox}
	}%
	\vskip0pt%
}

% Logo institucional (descomentar y ajustar ruta si tienes logo)
% \logo{\includegraphics[height=0.8cm]{logos/unsch_logo.png}}

% Fuentes personalizadas
\setbeamerfont{title}{size=\Large,series=\bfseries}
\setbeamerfont{frametitle}{size=\large,series=\bfseries}
\setbeamerfont{block title}{size=\normalsize,series=\bfseries}

% Espaciado entre items
\setlength{\itemsep}{0.5em}

% Configuración de listings (código)
\lstset{
	basicstyle=\ttfamily\footnotesize,
	keywordstyle=\color{blue}\bfseries,
	commentstyle=\color{green!50!black},
	stringstyle=\color{orange},
	numbers=left,
	numberstyle=\tiny\color{gray},
	stepnumber=1,
	numbersep=5pt,
	backgroundcolor=\color{gray!10},
	frame=single,
	breaklines=true,
	captionpos=b,
	showstringspaces=false
}

% Estilo para bloques de alerta/ejemplos
\setbeamertemplate{blocks}[rounded][shadow=true]

%%%%%%%%%%%%%%%%%%%%%%%%%%%%%%%%%%%%%%%%%%%%%%%%%%%%%%%%%%%%%%%%%
% INFORMACIÓN DEL DOCUMENTO
%%%%%%%%%%%%%%%%%%%%%%%%%%%%%%%%%%%%%%%%%%%%%%%%%%%%%%%%%%%%%%%%%

\title[IA para Universitarios]{Inteligencia Artificial Aplicada para Universitarios}
\subtitle{Ecosistema Completo de IA para Investigación, Redacción y Productividad Académica}
\author[E. Achalma]{Edison Achalma}
\institute[UNSCH]{
	Universidad Nacional de San Cristóbal de Huamanga\\
	\medskip
	\textit{achalmaedison@unsch.edu.pe}
}
\date{\today}

%%%%%%%%%%%%%%%%%%%%%%%%%%%%%%%%%%%%%%%%%%%%%%%%%%%%%%%%%%%%%%%%%
% COMANDOS PERSONALIZADOS
%%%%%%%%%%%%%%%%%%%%%%%%%%%%%%%%%%%%%%%%%%%%%%%%%%%%%%%%%%%%%%%%%

% Comando para crear slide de sección con estilo especial
\newcommand{\seccionslide}[1]{
	\begin{frame}[plain,noframenumbering]
		\vfill
		\centering
		\begin{beamercolorbox}[sep=12pt,center,shadow=true,rounded=true]{title}
			\usebeamerfont{title}\LARGE #1\par
		\end{beamercolorbox}
		\vfill
	\end{frame}
}

% Comando para ícono + texto (requiere fontawesome5)
\newcommand{\icono}[1]{\textcolor{azulUNSCH}{#1}}

% Comando para destacar herramientas
\newcommand{\herramienta}[1]{\textbf{\textcolor{naranjaUNSCH}{#1}}}

% Comando para URLs clickeables con estilo
\newcommand{\enlace}[2]{\href{#1}{\textcolor{cyan}{\underline{#2}}}}

%%%%%%%%%%%%%%%%%%%%%%%%%%%%%%%%%%%%%%%%%%%%%%%%%%%%%%%%%%%%%%%%%
% INICIO DEL DOCUMENTO
%%%%%%%%%%%%%%%%%%%%%%%%%%%%%%%%%%%%%%%%%%%%%%%%%%%%%%%%%%%%%%%%%

\begin{document}
	
	%%%%%%%%%%%%%%%%%%%%%%%%%%%%%%%%%%%%%%%%%%%%%%%%%%%%%%%%%%%%%%%%%
	% PORTADA
	%%%%%%%%%%%%%%%%%%%%%%%%%%%%%%%%%%%%%%%%%%%%%%%%%%%%%%%%%%%%%%%%%
	
	\begin{frame}
		\titlepage
		
		\vspace{-0.5cm}
		\centering
		\footnotesize
		\textit{Duración: 2 horas | 118 slides | 40+ herramientas de IA}
	\end{frame}
	
	%%%%%%%%%%%%%%%%%%%%%%%%%%%%%%%%%%%%%%%%%%%%%%%%%%%%%%%%%%%%%%%%%
	% TABLA DE CONTENIDO
	%%%%%%%%%%%%%%%%%%%%%%%%%%%%%%%%%%%%%%%%%%%%%%%%%%%%%%%%%%%%%%%%%
	
	\begin{frame}{Contenido de la Presentación}
		\tableofcontents
	\end{frame}
	
	%%%%%%%%%%%%%%%%%%%%%%%%%%%%%%%%%%%%%%%%%%%%%%%%%%%%%%%%%%%%%%%%%
	% SECCIÓN 1: APERTURA Y CONTEXTO
	%%%%%%%%%%%%%%%%%%%%%%%%%%%%%%%%%%%%%%%%%%%%%%%%%%%%%%%%%%%%%%%%%
	
	\section{Apertura y Contexto}
	
	% Slide de separación de sección
	\seccionslide{🚀 Apertura y Contexto}
	
	\begin{frame}{Bienvenida}
		\begin{center}
			{\LARGE \textbf{El Ecosistema de IA:}}\\[0.5cm]
			{\Large Más allá de ChatGPT}\\[1cm]
			
			\includegraphics[width=0.6\textwidth]{imagenes/ecosistema_ia.png}
			% Nota: Reemplazar con imagen real o crear con TikZ
			
			\vspace{0.5cm}
			{\large Tu caja de herramientas digital para la universidad}
		\end{center}
	\end{frame}
	
	\begin{frame}{¿Qué lograrás hoy?}
		
		Al finalizar esta sesión serás capaz de:
		
		\begin{itemize}
			\item \icono{\faCheck} Conocer \textbf{+40 herramientas} de IA especializadas
			
			\item \icono{\faRobot} Usar asistentes conversacionales avanzados:\\
			ChatGPT, Claude, Gemini, Grok, DeepSeek, Meta AI
			
			\item \icono{\faSearch} Investigar con IA académica:\\
			Perplexity, Consensus, Scite, Litmaps
			
			\item \icono{\faPalette} Crear contenido multimedia:\\
			Texto, imágenes, videos, audio, presentaciones
			
			\item \icono{\faCogs} Automatizar tareas repetitivas con IA
		\end{itemize}
		
		\vspace{0.3cm}
		\begin{alertblock}{Objetivo}
			Pasar de usuario casual a usuario estratégico de IA
		\end{alertblock}
		
	\end{frame}
	
	\begin{frame}{Encuesta Rápida Interactiva}
		
		\begin{center}
			{\large Levanta la mano si...}
		\end{center}
		
		\vspace{0.5cm}
		
		\begin{columns}[T]
			\begin{column}{0.48\textwidth}
				\begin{block}{Experiencia actual}
					\begin{itemize}
						\item ¿Has usado ChatGPT?
						\item ¿Conoces Claude o Gemini?
						\item ¿Has generado imágenes con IA?
						\item ¿Usas IA para estudiar?
					\end{itemize}
				\end{block}
			\end{column}
			
			\begin{column}{0.48\textwidth}
				\begin{exampleblock}{Meta de hoy}
					\centering
					Al final conocerás\\[0.3cm]
					{\Huge \textbf{40+}}\\[0.2cm]
					{\large herramientas especializadas}\\[0.5cm]
					
					\small (No solo ChatGPT)
				\end{exampleblock}
			\end{column}
		\end{columns}
		
	\end{frame}
	
	\begin{frame}{El Problema del Estudiante Moderno}
		
		\begin{center}
			\begin{tikzpicture}[node distance=2.5cm]
				% Nodo central
				\node[draw, circle, fill=naranjaUNSCH!20, minimum size=2.5cm, text width=2cm, align=center] (estudiante) 
				{\textbf{Estudiante\\2025}};
				
				% Problemas alrededor
				\node[draw, rectangle, fill=red!20, above left of=estudiante, text width=2.5cm] (sobrecarga) 
				{Sobrecarga de información};
				
				\node[draw, rectangle, fill=red!20, above right of=estudiante, text width=2.5cm] (tiempo) 
				{Tiempo limitado};
				
				\node[draw, rectangle, fill=red!20, below left of=estudiante, text width=2.5cm] (calidad) 
				{Contenido de calidad};
				
				\node[draw, rectangle, fill=red!20, below right of=estudiante, text width=2.5cm] (formatos) 
				{Múltiples formatos};
				
				% Flechas
				\draw[->,thick] (sobrecarga) -- (estudiante);
				\draw[->,thick] (tiempo) -- (estudiante);
				\draw[->,thick] (calidad) -- (estudiante);
				\draw[->,thick] (formatos) -- (estudiante);
			\end{tikzpicture}
		\end{center}
		
	\end{frame}
	
	\begin{frame}{La Solución: Ecosistema de IA Especializada}
		
		\begin{alertblock}{Principio fundamental}
			No existe \textbf{``la mejor IA''} $\rightarrow$ Cada herramienta tiene su caso de uso óptimo
		\end{alertblock}
		
		\vspace{0.5cm}
		
		\begin{columns}[T]
			\begin{column}{0.48\textwidth}
				\textbf{Conversación general:}
				\begin{itemize}
					\item ChatGPT
					\item Claude
					\item Gemini
				\end{itemize}
				
				\vspace{0.3cm}
				
				\textbf{Investigación académica:}
				\begin{itemize}
					\item Perplexity
					\item Consensus
					\item Scite
				\end{itemize}
			\end{column}
			
			\begin{column}{0.48\textwidth}
				\textbf{Creatividad visual:}
				\begin{itemize}
					\item DALL-E
					\item Midjourney
					\item Stable Diffusion
				\end{itemize}
				
				\vspace{0.3cm}
				
				\textbf{Productividad:}
				\begin{itemize}
					\item Fireflies.ai
					\item Otter.ai
					\item Gamma
				\end{itemize}
			\end{column}
		\end{columns}
		
	\end{frame}
	
	\begin{frame}{Mapa Mental del Ecosistema}
		
		\begin{center}
			\begin{tikzpicture}[
				mindmap,
				grow cyclic,
				every node/.style=concept,
				concept color=azulUNSCH!40,
				level 1/.style={level distance=4.5cm,sibling angle=90},
				level 2/.style={level distance=3cm,sibling angle=45}
				]
				
				\node[concept,scale=0.9]{\textbf{Tú}\\(Estudiante)}
				child[concept color=green!40] { 
					node[scale=0.7] {Asistentes\\Conversacionales}
				}
				child[concept color=blue!40] { 
					node[scale=0.7] {Investigación\\Académica}
				}
				child[concept color=orange!40] { 
					node[scale=0.7] {Creación de\\Contenido}
				}
				child[concept color=purple!40] { 
					node[scale=0.7] {Productividad\\Automatización}
				};
				
			\end{tikzpicture}
		\end{center}
		
		\vspace{0.3cm}
		\centering
		\footnotesize
		Cada rama representa un conjunto de herramientas especializadas
		
	\end{frame}
	
	\begin{frame}{Reglas de Oro para Hoy}
		
		\begin{enumerate}
			\item \icono{\faLaptop} \textbf{Trae laptop o celular} para probar en vivo
			
			\item \icono{\faClock} Cada herramienta = \textbf{2-3 minutos} de demostración
			
			\item \icono{\faQuestion} \textbf{Pregunta en cualquier momento}\\
			(No hay preguntas tontas, solo oportunidades de aprender)
			
			\item \icono{\faEdit} Tomaremos \textbf{notas colaborativas}\\
			(Enlace al documento compartido en el chat)
		\end{enumerate}
		
		\vspace{0.5cm}
		
		\begin{block}{Consejo}
			Marca con \icono{\faStar} las herramientas que más te interesen durante la presentación
		\end{block}
		
	\end{frame}
	
	\begin{frame}{Metodología de la Clase}
		
		\begin{table}
			\centering
			\small
			\begin{tabular}{clc}
				\toprule
				\textbf{Tiempo} & \textbf{Bloque} & \textbf{Herramientas} \\
				\midrule
				10 min & \icono{\faBook} Teoría: ¿Qué es la IA? & Conceptos \\
				24 min & \icono{\faRobot} Asistentes Conversacionales & 9 tools \\
				16 min & \icono{\faSearch} Investigación Académica & 8 tools \\
				\midrule
				10 min & \icono{\faCoffee} \textbf{BREAK} & Descanso \\
				\midrule
				26 min & \icono{\faPalette} Generación de Contenido & 12 tools \\
				19 min & \icono{\faCogs} Herramientas Especializadas & 8 tools \\
				7 min & \icono{\faBalanceScale} Ética + Cierre & Recursos \\
				\bottomrule
				\textbf{127 min} & \textbf{TOTAL} & \textbf{40+ tools} \\
				\bottomrule
			\end{tabular}
		\end{table}
		
		\vspace{0.3cm}
		\centering
		\footnotesize
		*Incluye 7 minutos de flexibilidad para interacciones y preguntas
		
	\end{frame}
	
	%%%%%%%%%%%%%%%%%%%%%%%%%%%%%%%%%%%%%%%%%%%%%%%%%%%%%%%%%%%%%%%%%
	% SECCIÓN 2: FUNDAMENTOS DE IA
	%%%%%%%%%%%%%%%%%%%%%%%%%%%%%%%%%%%%%%%%%%%%%%%%%%%%%%%%%%%%%%%%%
	
	\section{Fundamentos de IA para Universitarios}
	
	\seccionslide{Fundamentos de IA\\para Universitarios}
	
	\begin{frame}{¿Qué es la Inteligencia Artificial?}
		
		\begin{definition}[Definición simple]
			La \textbf{Inteligencia Artificial} es un programa de computadora que \textit{aprende patrones} a partir de ejemplos, sin ser explícitamente programado para cada tarea.
		\end{definition}
		
		\vspace{0.5cm}
		
		\begin{columns}[T]
			\begin{column}{0.48\textwidth}
				\textbf{Lo que SÍ es:}
				\begin{itemize}
					\item Reconocimiento de patrones
					\item Predicción basada en datos
					\item Imitación de comportamiento humano
				\end{itemize}
			\end{column}
			
			\begin{column}{0.48\textwidth}
				\textbf{Lo que NO es:}
				\begin{itemize}
					\item Pensamiento consciente
					\item Comprensión real
					\item Inteligencia como la humana
				\end{itemize}
			\end{column}
		\end{columns}
		
		\vspace{0.5cm}
		
		\begin{exampleblock}{Analogía}
			La IA es como un estudiante que ha leído millones de libros: puede responder muchas preguntas, pero no ``entiende'' como un humano.
		\end{exampleblock}
		
	\end{frame}
	
	\begin{frame}{Tipos de IA que Usarás en la Universidad}
		
		\begin{table}
			\centering
			\begin{tabular}{p{3.5cm}p{5cm}p{3cm}}
				\toprule
				\textbf{Tipo de IA} & \textbf{Qué hace} & \textbf{Ejemplo} \\
				\midrule
				\textbf{IA Generativa} & Crea contenido nuevo (texto, imagen, audio, video) & ChatGPT, DALL-E \\
				\addlinespace
				\textbf{IA de Búsqueda} & Encuentra y organiza información relevante & Perplexity, Consensus \\
				\addlinespace
				\textbf{IA de Análisis} & Procesa y extrae insights de datos/documentos & NotebookLM, Scite \\
				\bottomrule
			\end{tabular}
		\end{table}
		
		\vspace{0.5cm}
		
		\begin{alertblock}{Nota importante}
			En esta clase nos enfocaremos en las \textbf{2 primeras categorías}, que son las más útiles para estudiantes universitarios.
		\end{alertblock}
		
	\end{frame}
	
	\begin{frame}{¿Cómo Funciona un Modelo de Lenguaje?}
		
		\begin{center}
			\begin{tikzpicture}[node distance=1.5cm, auto]
				% Nodos
				\node[draw, rectangle, fill=blue!20, text width=3cm, align=center] (entrenamiento) 
				{1. Entrenamiento\\(Millones de textos)};
				
				\node[draw, rectangle, fill=green!20, text width=3cm, align=center, right=of entrenamiento] (aprendizaje) 
				{2. Aprende patrones\\(Relaciones entre palabras)};
				
				\node[draw, rectangle, fill=orange!20, text width=3cm, align=center, right=of aprendizaje] (prediccion) 
				{3. Predice\\(Siguiente palabra más probable)};
				
				% Flechas
				\draw[->,thick] (entrenamiento) -- (aprendizaje);
				\draw[->,thick] (aprendizaje) -- (prediccion);
			\end{tikzpicture}
		\end{center}
		
		\vspace{0.5cm}
		
		\begin{exampleblock}{Ejemplo concreto}
			Entrada: ``El gato está en el...''\\
			IA predice: ``tejado'' (80\%), ``árbol'' (15\%), ``coche'' (5\%)\\
			Resultado: ``El gato está en el tejado''
		\end{exampleblock}
		
		\vspace{0.3cm}
		
		\begin{alertblock}{Importante}
			La IA \textbf{no entiende} el concepto de ``gato'' o ``tejado''. Solo sabe que estadísticamente esas palabras aparecen juntas frecuentemente.
		\end{alertblock}
		
	\end{frame}
	
	\begin{frame}{Conceptos Clave de IA}
		
		\begin{description}[style=nextline]
			\item[\herramienta{Prompt}] 
			Instrucción que le das a la IA. Entre más detallado y claro, mejor la respuesta.
			
			\item[\herramienta{Token}] 
			Unidad de texto que procesa la IA ($\approx$ 4 caracteres). Los modelos tienen límites de tokens por conversación.
			
			\item[\herramienta{Contexto}] 
			Memoria de la conversación. La IA recuerda mensajes anteriores dentro de un límite.
			
			\item[\herramienta{Alucinación}] 
			Cuando la IA inventa información con confianza. \textbf{Siempre verifica datos críticos.}
		\end{description}
		
		\vspace{0.5cm}
		
		\begin{alertblock}{Ejemplo de alucinación}
			Pregunta: ``¿Quién ganó el Nobel de Literatura 2024?''\\
			IA (sin acceso a internet): ``John Smith ganó por su novela...'' \textcolor{red}{[INVENTADO]}
		\end{alertblock}
		
	\end{frame}
	
	\begin{frame}{Fortalezas de la IA para Estudiantes}
		
		\begin{columns}[T]
			\begin{column}{0.48\textwidth}
				\begin{block}{Lo que hace EXCELENTE:}
					\begin{itemize}
						\item Sintetizar información compleja
						\item Generar múltiples alternativas rápidamente
						\item Explicar conceptos en diferentes niveles
						\item Disponibilidad 24/7
						\item Traducir entre idiomas
					\end{itemize}
				\end{block}
			\end{column}
			
			\begin{column}{0.48\textwidth}
				\begin{exampleblock}{Casos de uso:}
					\begin{itemize}
						\item Resume 20 páginas en 2 párrafos
						\item Explica fotosíntesis como si tuvieras 10 años
						\item Genera 10 títulos para tu ensayo
						\item Traduce paper técnico del inglés
					\end{itemize}
				\end{exampleblock}
			\end{column}
		\end{columns}
		
		\vspace{0.5cm}
		
		\begin{center}
			\textit{``La IA es tu asistente más trabajador, pero tú sigues siendo el jefe''}
		\end{center}
		
	\end{frame}
	
	\begin{frame}{Limitaciones Críticas de la IA}
		
		\begin{alertblock}{PELIGROS REALES - SIEMPRE CONSIDERAR}
			\begin{enumerate}
				\item \textbf{Alucinaciones}: Puede inventar datos, referencias, estudios\\
				\textit{``El 87\% de las estadísticas en internet son inventadas''} - Abraham Lincoln \textcolor{red}{[FALSO]}
				
				\item \textbf{Fecha de corte}: ChatGPT 3.5 (Sep 2021), GPT-4 (Abr 2023)\\
				No sabe sobre eventos recientes sin acceso a internet
				
				\item \textbf{No accede a bases de datos de pago}:\\
				No puede leer papers en JSTOR, ScienceDirect, etc.
				
				\item \textbf{Sesgos en entrenamiento}: Refleja sesgos de sus datos\\
				Puede tener más info de EE.UU. que de Honduras
				
				\item \textbf{No reemplaza pensamiento crítico}:\\
				Tú decides qué es verdad, qué usar, cómo aplicar
			\end{enumerate}
		\end{alertblock}
		
	\end{frame}
	
	\begin{frame}{Ejemplo Real de Alucinación}
		
		\begin{exampleblock}{Caso documentado (2023)}
			\textbf{Pregunta a ChatGPT 3.5:}\\
			``Lista 5 papers sobre tratamiento de Alzheimer publicados en 2022''
			
			\vspace{0.3cm}
			
			\textbf{Respuesta de ChatGPT:}\\
			\small
			1. Smith et al. (2022) ``New drug reverses Alzheimer's'' - Nature\\
			2. Johnson (2022) ``Breakthrough in protein folding'' - Science\\
			3. García et al. (2022) ``Clinical trial results'' - NEJM\\
			[...]
			
			\vspace{0.3cm}
			
			\textbf{Realidad:}\\
			\textcolor{red}{NINGUNO de estos papers existe}. Fueron completamente inventados.
		\end{exampleblock}
		
		\vspace{0.3cm}
		
		\begin{alertblock}{Lección}
			SIEMPRE verifica información crítica (papers, datos, citas) en fuentes primarias o con herramientas que citan fuentes (Perplexity, Consensus).
		\end{alertblock}
		
	\end{frame}
	
	\begin{frame}{Gratuito vs Pago: ¿Qué Necesitas?}
		
		\begin{table}
			\centering
			\small
			\begin{tabular}{p{3cm}p{4.5cm}p{4.5cm}}
				\toprule
				\textbf{Aspecto} & \textbf{Versión Gratuita} & \textbf{Versión Pago} \\
				\midrule
				Modelo & GPT-3.5, Claude Instant & GPT-4, Claude Opus \\
				\addlinespace
				Velocidad & Más lenta en horas pico & Rápida siempre \\
				\addlinespace
				Límites & 40-50 msg/3h (varía) & Ilimitado o alto límite \\
				\addlinespace
				Funciones extra & Básicas & Plugins, DALL-E, análisis \\
				\addlinespace
				Precio & \$0/mes & \$20-30/mes \\
				\bottomrule
			\end{tabular}
		\end{table}
		
		\vspace{0.5cm}
		
		\begin{block}{Recomendación para estudiantes}
			\begin{itemize}
				\item \textbf{Empieza con versiones gratuitas} (suficiente para 80\% de tareas)
				\item Usa múltiples herramientas gratis (ChatGPT + Claude + Gemini)
				\item Paga solo si: (a) usas diariamente, (b) necesitas velocidad, (c) requieres funciones avanzadas
			\end{itemize}
		\end{block}
		
	\end{frame}
	
	\begin{frame}{¿Qué Herramienta para Qué Tarea?}
		
		\begin{table}
			\centering
			\footnotesize
			\begin{tabular}{p{4cm}p{7cm}}
				\toprule
				\textbf{Tarea} & \textbf{Herramienta Recomendada} \\
				\midrule
				Explicar conceptos & ChatGPT, Claude, Gemini \\
				\addlinespace
				Investigar papers académicos & Consensus, Scite, Connected Papers \\
				\addlinespace
				Búsqueda con fuentes verificables & Perplexity, Bing AI \\
				\addlinespace
				Crear presentaciones & Gamma, Tome \\
				\addlinespace
				Generar imágenes & DALL-E, Midjourney \\
				\addlinespace
				Transcribir clases & Otter.ai, Fireflies.ai \\
				\addlinespace
				Programar código & GitHub Copilot, Blackbox \\
				\addlinespace
				Analizar PDFs largos & ChatPDF, NotebookLM \\
				\bottomrule
			\end{tabular}
		\end{table}
		
		\vspace{0.3cm}
		\centering
		\footnotesize
		*Tabla detallada en documento colaborativo compartido
		
	\end{frame}
	
	\begin{frame}{Transición: De Teoría a Práctica}
		
		\begin{center}
			\begin{tikzpicture}
				\node[draw, circle, fill=blue!20, minimum size=3cm] (teoria) at (0,0) 
				{\begin{tabular}{c} \textbf{TEORÍA} \\ Fundamentos \\ de IA \end{tabular}};
				
				\node[draw, circle, fill=green!20, minimum size=3cm] (practica) at (6,0) 
				{\begin{tabular}{c} \textbf{PRÁCTICA} \\ 40+ \\ Herramientas \end{tabular}};
				
				\draw[->,ultra thick,red] (teoria) -- (practica) node[midway,above] {\Large AHORA};
			\end{tikzpicture}
		\end{center}
		
		\vspace{0.5cm}
		
		\begin{block}{Lo que viene:}
			\begin{itemize}
				\item \textbf{24 minutos}: Asistentes Conversacionales (9 herramientas)
				\item \textbf{16 minutos}: Investigación Académica (8 herramientas)
				\item Cada herramienta = 2 minutos de demo en vivo
				\item Toma nota de las que más te interesen
			\end{itemize}
		\end{block}
		
		\vspace{0.3cm}
		\centering
		{\large \textbf{Prepara tu laptop/celular para probar}}
		
	\end{frame}
	
	%%%%%%%%%%%%%%%%%%%%%%%%%%%%%%%%%%%%%%%%%%%%%%%%%%%%%%%%%%%%%%%%%
	% SECCIÓN 3: ASISTENTES CONVERSACIONALES
	%%%%%%%%%%%%%%%%%%%%%%%%%%%%%%%%%%%%%%%%%%%%%%%%%%%%%%%%%%%%%%%%%
	
	\section{Asistentes Conversacionales}
	
	\seccionslide{Bloque 1:\\Asistentes Conversacionales}
	
	% Nota: Por limitaciones de espacio, aquí presento la estructura de ALGUNOS slides clave
	% El documento completo incluiría todos los 118 slides siguiendo este formato
	
	\begin{frame}{Introducción: Asistentes Conversacionales}
		
		\begin{definition}[Asistentes Conversacionales]
			IAs de propósito general entrenadas para mantener conversaciones coherentes sobre cualquier tema. Son los ``cerebros generales'' del ecosistema de IA.
		\end{definition}
		
		\vspace{0.5cm}
		
		\begin{block}{Útiles para:}
			\begin{itemize}
				\item Explicaciones de conceptos complejos
				\item Brainstorming y generación de ideas
				\item Redacción y corrección de textos
				\item Resolución de problemas paso a paso
				\item Traducción y reformulación
			\end{itemize}
		\end{block}
		
		\vspace{0.3cm}
		
		\begin{alertblock}{Hoy veremos 6 modelos principales}
			ChatGPT • Claude • Gemini • DeepSeek • Grok • Meta AI
		\end{alertblock}
		
	\end{frame}
	
	\begin{frame}{\herramienta{ChatGPT} - OpenAI}
		
		\begin{columns}[T]
			\begin{column}{0.48\textwidth}
				\textbf{Versiones:}
				\begin{itemize}
					\item GPT-3.5 (gratis)
					\item GPT-4 (pago, \$20/mes)
					\item GPT-4o (más rápido)
				\end{itemize}
				
				\vspace{0.3cm}
				
				\textbf{Fortalezas:}
				\begin{itemize}
					\item Versatilidad
					\item Plugins disponibles
					\item DALL-E integrado (imágenes)
					\item Gran comunidad
				\end{itemize}
			\end{column}
			
			\begin{column}{0.48\textwidth}
				\textbf{Ideal para:}
				\begin{itemize}
					\item Redacción de ensayos
					\item Explicaciones de conceptos
					\item Programación
					\item Traducción
				\end{itemize}
				
				\vspace{0.3cm}
				
				\textbf{Límites gratuitos:}
				\begin{itemize}
					\item Modelo 3.5: Ilimitado
					\item Modelo 4: ~40 msg/3h
				\end{itemize}
			\end{column}
		\end{columns}
		
		\vspace{0.5cm}
		
		\begin{exampleblock}{Demo en vivo (2 min)}
			\enlace{https://chat.openai.com}{chat.openai.com}\\
			Tarea: ``Explica la fotosíntesis como si tuviera 10 años''\\
			Luego: ``Ahora explícala para ensayo universitario de Biología''
		\end{exampleblock}
		
	\end{frame}
	
	% [CONTINÚA CON CLAUDE, GEMINI, DEEPSEEK, GROK, META AI...]
	% Por brevedad, omito los slides intermedios pero seguirían el mismo formato
	
	%%%%%%%%%%%%%%%%%%%%%%%%%%%%%%%%%%%%%%%%%%%%%%%%%%%%%%%%%%%%%%%%%
	% SLIDE DE EJERCICIO PRÁCTICO
	%%%%%%%%%%%%%%%%%%%%%%%%%%%%%%%%%%%%%%%%%%%%%%%%%%%%%%%%%%%%%%%%%
	
	\begin{frame}{Ejercicio Práctico 1}
		
		\begin{alertblock}{Instrucciones (5 minutos)}
			\begin{enumerate}
				\item Elige 2 asistentes diferentes (ej: ChatGPT y Claude)
				\item Haz la MISMA pregunta en ambos
				\item Compara las respuestas
			\end{enumerate}
		\end{alertblock}
		
		\vspace{0.3cm}
		
		\begin{exampleblock}{Pregunta sugerida}
			``Explica en 200 palabras cómo funciona la fotosíntesis para un ensayo de Biología universitaria''
		\end{exampleblock}
		
		\vspace{0.5cm}
		
		\begin{block}{Qué observar:}
			\begin{itemize}
				\item ¿Cuál dio respuesta más clara?
				\item ¿Cuál usó mejor vocabulario técnico?
				\item ¿Cuál organizó mejor la información?
				\item ¿Cuál preferirías usar en el futuro?
			\end{itemize}
		\end{block}
		
		\vspace{0.3cm}
		\centering
		{\large Compartiremos observaciones en 5 minutos}
		
	\end{frame}
	
	%%%%%%%%%%%%%%%%%%%%%%%%%%%%%%%%%%%%%%%%%%%%%%%%%%%%%%%%%%%%%%%%%
	% SLIDES DE RESUMEN DE BLOQUES (EJEMPLO)
	%%%%%%%%%%%%%%%%%%%%%%%%%%%%%%%%%%%%%%%%%%%%%%%%%%%%%%%%%%%%%%%%%
	
	\begin{frame}{Resumen: Asistentes Conversacionales}
		
		\begin{table}
			\centering
			\footnotesize
			\begin{tabular}{lccc}
				\toprule
				\textbf{Herramienta} & \textbf{Mejor para} & \textbf{Gratis} & \textbf{Link} \\
				\midrule
				ChatGPT & Versatilidad & & chat.openai.com \\
				Claude & Textos largos & & claude.ai \\
				Gemini & Info actualizada & & gemini.google.com \\
				DeepSeek & Matemáticas & & chat.deepseek.com \\
				Grok & Tendencias X & & x.ai \\
				Meta AI & WhatsApp/IG & & En apps Meta \\
				\bottomrule
			\end{tabular}
		\end{table}
		
		\vspace{0.5cm}
		
		\begin{block}{Qué aprendiste:}
			\begin{itemize}
				\item Conociste 6 asistentes conversacionales principales
				\item Probaste al menos 2 en vivo
				\item Entendiste diferencias: velocidad, límites, especialización
				\item Sabes cuál usar según tu tarea
			\end{itemize}
		\end{block}
		
		\vspace{0.3cm}
		\centering
		{\large \textbf{Siguiente:} Buscadores con IA para Investigación}
		
	\end{frame}
	
	%%%%%%%%%%%%%%%%%%%%%%%%%%%%%%%%%%%%%%%%%%%%%%%%%%%%%%%%%%%%%%%%%
	% SECCIONES SIGUIENTES (ESTRUCTURA SIMILAR)
	%%%%%%%%%%%%%%%%%%%%%%%%%%%%%%%%%%%%%%%%%%%%%%%%%%%%%%%%%%%%%%%%%
	
	% NOTA: Para no exceder límites de respuesta, presento solo la estructura
	% El documento completo incluiría todos los bloques con este nivel de detalle
	
	\section{Buscadores con IA para Investigación}
	\seccionslide{Bloque 2:\\Buscadores con IA\\para Investigación}
	
	% [Aquí irían slides 43-58 con el mismo formato]
	% - Consensus (demo)
	% - Scite (demo)
	% - Litmaps (demo)
	% - Connected Papers (demo)
	% - NotebookLM (demo)
	% - SciSpace (demo)
	% - ChatPDF (demo)
	% - Ejercicio práctico 2
	% - Resumen del bloque
	
	%%%%%%%%%%%%%%%%%%%%%%%%%%%%%%%%%%%%%%%%%%%%%%%%%%%%%%%%%%%%%%%%%
	% BREAK
	%%%%%%%%%%%%%%%%%%%%%%%%%%%%%%%%%%%%%%%%%%%%%%%%%%%%%%%%%%%%%%%%%
	
	\section{Break}
	
	\begin{frame}[plain]
		\vfill
		\centering
		{\Huge \textbf{BREAK}}\\[1cm]
		{\LARGE 10 minutos de descanso}\\[1.5cm]
		
		\begin{block}{Desafío opcional}
			Busca 1 paper de tu área en \textbf{Consensus}\\
			Verifica su confiabilidad en \textbf{Scite}\\
			Descubre papers relacionados en \textbf{Connected Papers}
		\end{block}
		
		\vspace{1cm}
		{\large Regresamos para: Generación de Contenido Visual}
		\vfill
	\end{frame}
	
	%%%%%%%%%%%%%%%%%%%%%%%%%%%%%%%%%%%%%%%%%%%%%%%%%%%%%%%%%%%%%%%%%
	% SECCIONES RESTANTES
	%%%%%%%%%%%%%%%%%%%%%%%%%%%%%%%%%%%%%%%%%%%%%%%%%%%%%%%%%%%%%%%%%
	
	\section{Generación de Contenido}
	\seccionslide{Bloque 3:\\Generación de Contenido}
	
	% [Slides 60-85: Tome, Gamma, DALL-E, Midjourney, Fliki, etc.]
	
	\section{Herramientas Especializadas}
	\seccionslide{Bloque 4:\\Herramientas Especializadas}
	
	% [Slides 86-104: Copilot, Blackbox, Fireflies, Otter, Directorios]
	
	\section{Ingeniería de Prompts y Ética}
	\seccionslide{Ingeniería de Prompts\\y Uso Ético}
	
	\begin{frame}{Anatomía de un Prompt Efectivo}
		
		\begin{block}{Estructura de 5 componentes}
			\begin{enumerate}
				\item \textbf{Rol}: ``Actúa como [experto en X]''
				\item \textbf{Contexto}: ``Para [audiencia/propósito]''
				\item \textbf{Tarea}: ``Genera/Explica/Analiza [Y]''
				\item \textbf{Formato}: ``En forma de [lista/tabla/párrafos]''
				\item \textbf{Restricciones}: ``Máximo X palabras, nivel Z''
			\end{enumerate}
		\end{block}
		
		\vspace{0.3cm}
		
		\begin{exampleblock}{Ejemplo completo}
			\small
			\textbf{Prompt malo:} ``Háblame de economía''\\[0.3cm]
			
			\textbf{Prompt bueno:}\\
			``Actúa como profesor de Economía. Para estudiante de 1er año de Administración en Honduras, explica el concepto de inflación en 3 párrafos, usando ejemplos del mercado local hondureño. Formato: párrafos académicos de 100 palabras cada uno.''
		\end{exampleblock}
		
	\end{frame}
	
	\begin{frame}{Decálogo Ético del Uso de IA}
		
		\begin{enumerate}
			\item \textbf{Respeto a la dignidad humana}: No uses IA para discriminar
			
			\item \textbf{Transparencia}: Declara cuando usas IA en trabajos formales
			
			\item \textbf{Veracidad}: Verifica información crítica siempre
			
			\item \textbf{Protección de datos}: No compartas info personal/sensible
			
			\item \textbf{Autonomía}: La IA ayuda, \textbf{tú decides}
			
			\item \textbf{No maleficencia}: No uses IA para hacer daño
			
			\item \textbf{Equidad e inclusión}: Cuestiona sesgos en respuestas
			
			\item \textbf{Responsabilidad}: Eres responsable del output final
			
			\item \textbf{Bien común}: Usa IA para mejorar, no para engañar
			
			\item \textbf{Derechos de autor}: Cita fuentes originales, no la IA
		\end{enumerate}
		
	\end{frame}
	
	\begin{frame}{¿Cuándo es Trampa vs Uso Legítimo?}
		
		\begin{columns}[T]
			\begin{column}{0.48\textwidth}
				\begin{alertblock}{TRAMPA}
					\begin{itemize}
						\item Copiar respuesta IA sin leer
						\item Presentar como 100\% tuyo
						\item Usar en exámenes prohibido
						\item Inventar referencias de IA
					\end{itemize}
				\end{alertblock}
			\end{column}
			
			\begin{column}{0.48\textwidth}
				\begin{exampleblock}{USO LEGÍTIMO}
					\begin{itemize}
						\item IA genera esquema → Investigas → Escribes
						\item Revisar gramática/estilo
						\item Explicar conceptos difíciles
						\item Practicar con exámenes IA
					\end{itemize}
				\end{exampleblock}
			\end{column}
		\end{columns}
		
		\vspace{0.5cm}
		
		\begin{block}{Regla de oro}
			\centering
			{\large La IA es tu \textbf{tutor}, no tu \textbf{sustituto}}\\[0.3cm]
			
			Tu cerebro debe \textbf{trabajar}, no solo \textbf{copiar}
		\end{block}
		
	\end{frame}
	
	%%%%%%%%%%%%%%%%%%%%%%%%%%%%%%%%%%%%%%%%%%%%%%%%%%%%%%%%%%%%%%%%%
	% SECCIÓN FINAL: CIERRE Y RECURSOS
	%%%%%%%%%%%%%%%%%%%%%%%%%%%%%%%%%%%%%%%%%%%%%%%%%%%%%%%%%%%%%%%%%
	
	\section{Cierre y Recursos}
	
	\begin{frame}{Recapitulación: ¿Qué Aprendimos Hoy?}
		
		\begin{columns}[T]
			\begin{column}{0.48\textwidth}
				\textbf{Asistentes (6):}
				\begin{itemize}
					\item ChatGPT
					\item Claude
					\item Gemini
					\item DeepSeek
					\item Grok
					\item Meta AI
				\end{itemize}
				
				\vspace{0.3cm}
				
				\textbf{Investigación (8):}
				\begin{itemize}
					\item Consensus
					\item Scite
					\item Litmaps
					\item Connected Papers
					\item NotebookLM
					\item SciSpace
					\item ChatPDF
					\item Perplexity
				\end{itemize}
			\end{column}
			
			\begin{column}{0.48\textwidth}
				\textbf{Contenido (12):}
				\begin{itemize}
					\item Tome, Gamma
					\item DALL-E, Midjourney
					\item Fliki, Runway
					\item Enhance Speech
					\item [...]
				\end{itemize}
				
				\vspace{0.3cm}
				
				\textbf{Especialización (8):}
				\begin{itemize}
					\item Copilot, Blackbox
					\item Fireflies, Otter
					\item Toolify, AIFINDY
					\item [...]
				\end{itemize}
			\end{column}
		\end{columns}
		
		\vspace{0.5cm}
		\centering
		{\LARGE \textbf{Total: 40+ herramientas}}
		
	\end{frame}
	
	\begin{frame}{Tu Kit Esencial de IA (Top 10)}
		
		\begin{table}
			\centering
			\begin{tabular}{clp{5cm}}
				\toprule
				\# & \textbf{Herramienta} & \textbf{Para qué} \\
				\midrule
				1 & ChatGPT & Uso general, versatilidad \\
				2 & Claude & Analizar textos largos (papers, docs) \\
				3 & Perplexity & Búsqueda con fuentes verificables \\
				4 & Consensus & Investigación académica científica \\
				5 & NotebookLM & Organizar investigación de tesis \\
				6 & Gamma & Crear presentaciones rápidas \\
				7 & DALL-E & Generar imágenes para proyectos \\
				8 & Otter.ai & Transcribir clases y reuniones \\
				9 & Copilot & Programar (si estudias CS/Ingeniería) \\
				10 & Toolify & Descubrir nuevas herramientas \\
				\bottomrule
			\end{tabular}
		\end{table}
		
		\vspace{0.3cm}
		\centering
		\footnotesize
		Descarga lista completa en: \enlace{https://bit.ly/ia-universitarios}{bit.ly/ia-universitarios}
		
	\end{frame}
	
	\begin{frame}{Plan de Acción: Próximos Pasos}
		
		\begin{block}{HOY (antes de dormir)}
			\begin{itemize}
				\item Crea cuenta gratuita en: \textbf{ChatGPT} y \textbf{Perplexity}
				\item Marca 3 herramientas que más te interesaron
			\end{itemize}
		\end{block}
		
		\begin{block}{ESTA SEMANA}
			\begin{itemize}
				\item Prueba 1 herramienta nueva por día
				\item Compara resultados entre diferentes IA
				\item Practica escribir prompts efectivos
			\end{itemize}
		\end{block}
		
		\begin{block}{ESTE MES}
			\begin{itemize}
				\item Aplica IA a 1 tarea real de clase
				\item Documenta qué funcionó y qué no
				\item Comparte aprendizajes con compañeros
			\end{itemize}
		\end{block}
		
		\begin{alertblock}{PERMANENTE}
			Sigue aprendiendo - La IA evoluciona \textbf{cada semana}
		\end{alertblock}
		
	\end{frame}
	
	\begin{frame}{Recursos para Seguir Aprendiendo}
		
		\begin{description}[style=nextline]
			\item[\icono{\faNewspaper} Newsletter]
			\textbf{The Rundown AI} - Resumen diario de noticias de IA\\
			\enlace{https://therundown.ai}{therundown.ai}
			
			\item[\icono{\faYoutube} YouTube]
			\textbf{Matt Wolfe} - Reviews semanales de herramientas\\
			\enlace{https://youtube.com/@mreflow}{youtube.com/@mreflow}
			
			\item[\icono{\faReddit} Comunidad]
			\textbf{r/ChatGPT} y \textbf{r/ArtificialIntelligence}\\
			Preguntas, trucos, descubrimientos
			
			\item[\icono{\faGraduationCap} Curso gratis]
			\textbf{``Generative AI for Everyone''} - Andrew Ng (Coursera)\\
			\enlace{https://coursera.org/learn/generative-ai-for-everyone}{coursera.org/gen-ai}
			
			\item[\icono{\faBook} Documentación]
			Documento colaborativo de esta clase (links, prompts, notas)\\
			\enlace{https://bit.ly/ia-unsch-2025}{bit.ly/ia-unsch-2025}
		\end{description}
		
	\end{frame}
	
	\begin{frame}[standout]
		\vfill
		\centering
		{\Huge \textbf{¡Gracias!}}\\[1cm]
		
		{\Large ¿Preguntas?}\\[1.5cm]
		
		\begin{columns}[c]
			\begin{column}{0.48\textwidth}
				\centering
				\icono{\faEnvelope} \textbf{Contacto:}\\
				achalmaedison@unsch.edu.pe
			\end{column}
			
			\begin{column}{0.48\textwidth}
				\centering
				\icono{\faLink} \textbf{Materiales:}\\
				\enlace{https://bit.ly/ia-unsch-slides}{bit.ly/ia-unsch-slides}
			\end{column}
		\end{columns}
		
		\vspace{2cm}
		
		{\large \textit{``La IA no reemplaza tu inteligencia,}}\\
		{\large \textit{amplifica tu potencial''}} \icono{\faRocket}
		
		\vfill
	\end{frame}
	
	%%%%%%%%%%%%%%%%%%%%%%%%%%%%%%%%%%%%%%%%%%%%%%%%%%%%%%%%%%%%%%%%%
	% FIN DEL DOCUMENTO
	%%%%%%%%%%%%%%%%%%%%%%%%%%%%%%%%%%%%%%%%%%%%%%%%%%%%%%%%%%%%%%%%%
	
\end{document}